% ============================================================================
% SECCIÓN §2.3: COMPUTACIÓN DISTRIBUIDA (Edge vs Cloud)
% Reformulada: reducción de extensión según indicación del revisor
% ============================================================================

\subsection{Procesamiento Distribuido: Computación en la Nube vs en el Borde}
\label{sec:edge-vs-cloud}

El despliegue masivo de dispositivos IoT plantea una decisión arquitectónica fundamental: centralizar el procesamiento en centros de datos en la nube (\textit{cloud}), o distribuirlo en pasarelas de borde (\textit{edge gateways}) cercanas a los dispositivos. La literatura reciente~\cite{andriuloEdgeComputingCloud2024,minhEdgeComputingIoTEnabled2022} coincide en que esta dicotomía no es binaria, sino un espectro de arquitecturas híbridas que equilibran latencia, ancho de banda, escalabilidad y costo según los requisitos de cada aplicación.

\subsubsection{Computación en la Nube (\textit{Cloud Computing})}

El paradigma de computación en la nube concentra capacidad de cómputo y almacenamiento en centros de datos remotos (AWS, Azure, GCP) con arquitectura virtualizada multiusuario~\cite{andriuloEdgeComputingCloud2024}. Sus principales ventajas incluyen escalabilidad horizontal prácticamente ilimitada, servicios gestionados que reducen la carga operacional, acceso a capacidades avanzadas de aprendizaje automático (\textit{machine learning}), y alta disponibilidad con replicación geo-distribuida (SLA 99,95-99,99\%).

Sin embargo, para aplicaciones de energía inteligente (\textit{Smart Energy}), presenta limitaciones críticas:

\begin{itemize}
\item \textbf{Latencia inaceptable}: Andriulo et al.~\cite{andriuloEdgeComputingCloud2024} reportan latencias típicas de 50-200~ms en arquitecturas centradas en la nube, incompatibles con los requisitos IEEE~2030.5 de control DER (<100~ms) y monitoreo de calidad de energía (<10~ms)~\cite{IEEERecommendedPractice}.
\item \textbf{Consumo excesivo de ancho de banda}: La transmisión continua de formas de onda sin filtrado local puede saturar enlaces WAN en despliegues de mediana escala.
\item \textbf{Dependencia de conectividad WAN}: Durante interrupciones, los sistemas solo-nube (\textit{cloud-only}) quedan completamente inoperativos.
\item \textbf{Costos OPEX lineales}: El modelo de pago por uso escala linealmente con el volumen de datos, resultando en costos elevados para telemetría de alta frecuencia.
\end{itemize}

\subsubsection{Computación en el Borde (\textit{Edge Computing})}

La computación en el borde distribuye el procesamiento hacia pasarelas locales, procesando datos cerca de su origen antes de transmitir hacia la nube solo agregados, alarmas o eventos relevantes~\cite{andriuloEdgeComputingCloud2024,minhEdgeComputingIoTEnabled2022}. Sus ventajas principales son:

\begin{itemize}
\item \textbf{Latencia ultra-baja}: Típicamente 1-10~ms para procesamiento local sin ida y vuelta WAN~\cite{andriuloEdgeComputingCloud2024}, habilitando control en tiempo real (detección de huecos de tensión, respuesta a la demanda).
\item \textbf{Reducción de ancho de banda 70-90\%}: Mediante filtrado, agregación temporal y compresión local antes del envío selectivo a la nube~\cite{andriuloEdgeComputingCloud2024}.
\item \textbf{Operación autónoma}: Funcionalidad ininterrumpida durante desconexiones WAN con almacenamiento local y motor de reglas.
\item \textbf{Privacidad por diseño}: Procesamiento local de datos granulares de consumo sin transmisión a la nube pública.
\end{itemize}

Sus limitaciones incluyen recursos computacionales restringidos (impracticable para modelos pesados de aprendizaje profundo), complejidad de escalamiento horizontal (requiere instalación física) y riesgo de punto único de fallo local.

\subsubsection{Arquitectura Híbrida: Solución Adoptada}

La literatura~\cite{andriuloEdgeComputingCloud2024,bonomiFogComputingItsRole2012} establece que las arquitecturas híbridas que combinan computación en el borde y en la nube son la solución óptima para IoT industrial. El principio rector consiste en ejecutar cada tarea en el nivel más bajo (cercano a los datos) que satisfaga sus requisitos, escalando hacia la nube solo cuando sea necesario. La Tabla~\ref{tab:edge-cloud-workload-distribution} presenta la distribución de tareas adoptada en esta tesis.

\begin{table}[H]
\centering
\caption{Distribución de tareas en la arquitectura híbrida para energía inteligente}
\label{tab:edge-cloud-workload-distribution}
\small
\resizebox{\textwidth}{!}{%
\begin{tabular}{|l|l|l|l|}
\hline
\rowcolor{gray!20}
\textbf{Tarea} & \textbf{Latencia req.} & \textbf{Ubicación} & \textbf{Justificación} \\
\hline
Control DER (desconexión) & <100~ms & Borde & Latencia nube 50-200~ms insuficiente \\
\hline
Monitoreo calidad de energía & <10~ms & Borde & Detección de huecos <1 ciclo AC \\
\hline
Compresión de formas de onda & <50~ms & Borde & Reduce ancho de banda 90\% \\
\hline
Detección de anomalías & <1~s & Borde & Alertas inmediatas sin dependencia WAN \\
\hline
Paneles de control locales & <2~s & Borde & Operación durante interrupciones WAN \\
\hline
Pronóstico de consumo & 1-10~min & Nube & Requiere modelos ML pesados + históricos \\
\hline
Analítica agregada multi-sitio & 10-60~min & Nube & Agregación de millones de medidores \\
\hline
Almacenamiento largo plazo & N/A & Nube & Políticas de ciclo de vida costo-efectivas \\
\hline
\end{tabular}}
\end{table}

La Tabla~\ref{tab:edge-cloud-latency-comparison} consolida las métricas cuantitativas de latencia que fundamentan esta distribución.

\begin{table}[H]
\centering
\caption{Comparación cuantitativa de latencia: computación en el borde vs en la nube}
\label{tab:edge-cloud-latency-comparison}
\small
\resizebox{\textwidth}{!}{%
\begin{tabular}{|l|l|l|l|}
\hline
\rowcolor{gray!20}
\textbf{Componente} & \textbf{Borde} & \textbf{Nube} & \textbf{Referencia} \\
\hline
Enlace ascendente dispositivo $\rightarrow$ procesador & 1-5~ms (Thread/HaLow local) & 20-80~ms (LTE WAN) & \cite{andriuloEdgeComputingCloud2024} \\
\hline
Tiempo de procesamiento & 2-5~ms (pasarela local) & 10-30~ms (ingesta en nube) & \cite{andriuloEdgeComputingCloud2024} \\
\hline
Enrutamiento WAN & 0~ms & 10-50~ms (Internet) & \cite{andriuloEdgeComputingCloud2024} \\
\hline
Enlace descendente de respuesta & 0~ms (local) & 20-80~ms (LTE WAN) & \cite{andriuloEdgeComputingCloud2024} \\
\hline
\textbf{Latencia TOTAL} & \textcolor{green}{\textbf{1-10~ms}} & \textcolor{red}{\textbf{50-200~ms}} & \cite{andriuloEdgeComputingCloud2024} \\
\hline
Requisito IEEE 2030.5 DER & \textcolor{green}{\checkmark{} Cumple <100~ms} & \textcolor{red}{\xmark{} Excede >100~ms} & \cite{IEEERecommendedPractice} \\
\hline
Calidad de energía (<10~ms) & \textcolor{green}{\checkmark{} Viable} & \textcolor{red}{\xmark{} No viable} & \cite{minhEdgeComputingIoTEnabled2022} \\
\hline
\end{tabular}}
\end{table}

En síntesis, la computación en el borde reduce la latencia de extremo a extremo por un factor de 5-20$\times$ y el ancho de banda WAN en 70-90\%. Estas mejoras no son incrementales sino \textbf{cualitativas}, habilitando categorías completas de aplicaciones de control en tiempo real imposibles con arquitecturas solo-nube. La arquitectura presentada en los Capítulos~3 y~4 implementa este modelo híbrido: procesamiento local en pasarelas Raspberry Pi con ThingsBoard Edge, TimescaleDB y Kafka~\cite{krepsKafkaDistributedMessaging2011} para tareas de tiempo real (<10~ms), y sincronización selectiva con la nube para analítica avanzada y almacenamiento a largo plazo.

% ============================================================================
% FIN SECCIÓN §2.3
% ============================================================================
